\documentclass[letterpaper,11pt]{article}

\usepackage[spanish,es-tabla]{babel}
\usepackage[utf8]{inputenc}
\usepackage[T1]{fontenc}
\usepackage{lmodern}
\usepackage[top=1.6cm,bottom=1.8cm,left=2.2cm,right=2.2cm,headheight=28pt,includehead,headsep=0.5cm]{geometry}
\usepackage{graphicx}
\usepackage{xcolor}
\usepackage{booktabs}
\usepackage{array}
\usepackage{float}
\usepackage{enumitem}
\usepackage{fancyhdr}
\usepackage{hyperref}
\usepackage{amsmath}
\usepackage{amssymb}
\usepackage{tikz}
\usepackage{caption}

\definecolor{verde}{RGB}{0,128,0}
\definecolor{azulmarino}{RGB}{0,0,128}
\definecolor{gris}{RGB}{95,95,95}

\hypersetup{colorlinks=true,linkcolor=azulmarino,urlcolor=azulmarino}

\newcommand{\universidad}{Universidad Técnica Metropolitana}
\newcommand{\carrera}{Ingeniería Civil en Ciencia de Datos}
\newcommand{\codcurso}{INFB8090}
\newcommand{\nombrecurso}{Computación Paralela y Distribuida}
\newcommand{\fecha}{Mayo de 2026}
\newcommand{\profesor}{Michael Gabriel Miranda Sandoval}

\fancyhf{}
\fancyheadoffset[L]{0.4cm}
\fancyhead[L]{
\begin{tikzpicture}[remember picture,overlay]
    \fill[verde] (current page.north west) rectangle ++(\paperwidth/3,-0.18cm);
    \fill[azulmarino] ([xshift=\paperwidth/3]current page.north west) rectangle ++(\paperwidth*2/3,-0.18cm);
    \fill[azulmarino] (current page.south west) rectangle ++(\paperwidth*0.75,0.18cm);
    \fill[verde] ([xshift=0.75\paperwidth]current page.south west) rectangle ++(\paperwidth*0.25,0.18cm);
\end{tikzpicture}
\rmfamily \fontsize{9}{11}\selectfont
{\textcolor{verde}{\codcurso} -
\textcolor{azulmarino}{\nombrecurso}\\
\textcolor{gris}{\carrera}}
}
\fancyfoot[C]{
    \fontsize{7pt}{7pt}\selectfont \color{gris}
    \universidad\ - Proyecto Integrador semestral \hfill \thepage
}
\renewcommand{\headrulewidth}{0pt}
\renewcommand{\footrulewidth}{0.05pt}

\setlength{\parskip}{0.25em}

\begin{document}
\pagestyle{fancy}

\begin{center}
{\Large\bfseries Proyecto Integrador semestral\par}
\vspace{0.15cm}
{\large\itshape Detección de anomalías de comportamiento en trayectorias ADS-B de OpenSky Network a escala anual mediante un pipeline híbrido Dask + OpenMP\par}
\vspace{0.25cm}
\begin{tabular}{rl}
\textbf{Integrantes:} & Welinton Barrera Mondaca, Joaquin Araya, Juan Toledo \\
\textbf{Profesor:} & \profesor \\
\textbf{Curso / Sección:} & \codcurso\ - \nombrecurso\ / 412 \\
\textbf{Fecha:} & \fecha \\
\end{tabular}
\end{center}
\vspace{0.15cm}
\hrule
\vspace{0.25cm}

\section*{1. Problema delimitado y pertinencia para ciencia de datos}

La hipótesis aplicada es que dentro del tráfico aéreo comercial existen trayectorias con desviaciones sistemáticas respecto del patrón esperado para una ruta dada: maniobras evasivas, aproximaciones interrumpidas (\textit{go-around}), holdings prolongados, descenso anómalo, cambios de rumbo no compatibles con condiciones meteorológicas, o ejecución de operaciones no comerciales (medevac, vuelos estatales, ferry sin pasajeros). Detectar y rankear estas anomalías a escala anual habilita análisis de seguridad operacional, auditoría de eficiencia de rutas y caracterización de tráfico no rutinario sobre regiones específicas.

El problema es de ciencia de datos y no de ingeniería plana porque exige extracción de \textit{features} comportamentales no triviales (curvatura por segmento, distancia DTW contra una ruta canónica, energía cinética estimada, dispersión angular), selección de un modelo de \textit{outlier scoring} robusto a clases desbalanceadas y validación con \textit{ground truth} parcial publicado por entes regulatorios. La escala es la que justifica el componente paralelo: una solución secuencial sobre un año de datos OpenSky no termina en tiempo razonable en una sola máquina, y un \textit{downsampling} arbitrario destruye precisamente la fenomenología que se quiere detectar.

\section*{2. Datos: volumen, tipo, estructura y origen}

Se trabajará sobre la tabla \texttt{state\_vectors\_data4} del histórico público de OpenSky Network (Trino/Impala). El esquema usado por etapa de \textit{feature engineering} incluye al menos \texttt{icao24} (identificador 24-bit del transponder), \texttt{time}, \texttt{lat}, \texttt{lon}, \texttt{baroaltitude}, \texttt{velocity}, \texttt{heading}, \texttt{vertrate}, \texttt{onground} y \texttt{callsign}. La frecuencia de muestreo efectiva es cercana a 1 Hz por aeronave cuando hay cobertura. El volumen comprimido en formato columnar es del orden de \textbf{50--100 GB por mes} en AVRO, por lo que un año supera \textbf{1 TB comprimido} ($\approx$ 6--10 TB descomprimidos). El particionamiento natural en origen es por \texttt{(hour, day)}. Origen y licenciamiento: OpenSky Network, uso académico permitido bajo registro institucional.

Como referencia parcial de \textit{ground truth} se usarán reportes públicos de incidentes (NTSB, ASN) cruzados por \texttt{callsign} y ventana temporal, y un conjunto de \textbf{anomalías sintéticas} inyectadas (perturbaciones controladas sobre trayectorias normales) para medir \textit{recall} cuando el \textit{ground truth} real es escaso.

\section*{3. Hipótesis de paralelización y distribución}

El pipeline no tiene una sola carga dominante: cada etapa exhibe un perfil distinto y exige una estrategia distinta. La hipótesis es que mezclar paralelismo de datos a nivel de archivo, distribución por \texttt{groupby} de Dask y paralelismo de memoria compartida con OpenMP en los \textit{kernels} numéricos densos maximiza el aprovechamiento del cluster sin que el GIL bloquee las fases CPU-bound de Python.

\begin{table}[H]
\centering
\caption{Mapa etapa $\to$ tipo de carga $\to$ estrategia.}
\begin{tabular}{p{0.20\linewidth} p{0.20\linewidth} p{0.55\linewidth}}
\toprule
Etapa & Tipo de carga & Estrategia y justificación \\
\midrule
Ingesta AVRO/Parquet & I/O-bound + descompresión & ThreadPool de Dask: lecturas concurrentes liberan GIL en C (pyarrow). Solapa I/O con descompresión. \\
\textit{Groupby} por \texttt{icao24} & Shuffle distribuido & Dask Distributed con particiones por hash de \texttt{icao24}: evita que un solo proceso construya todas las trayectorias en RAM. \\
Extracción de \textit{features} & CPU-bound denso & ProcessPool + \textit{kernels} en C++ con \texttt{\#pragma omp parallel for} para DTW, distancia haversine y dispersión angular sobre trayectoria. Evita GIL y mantiene localidad. \\
\textit{Scoring} de anomalías & CPU-bound, paralelo por partición & IsolationForest entrenado distribuido, inferencia \textit{embarrassingly parallel} sobre particiones de Dask. \\
\bottomrule
\end{tabular}
\end{table}

\section*{4. Estrategia técnica preliminar}

El pipeline se diseña en cuatro etapas con materialización intermedia en Parquet para permitir reinicio sin reprocesar, y un \textit{scheduler} Dask con \textit{workers} heterogéneos (más memoria para el \textit{shuffle}, más CPU para \textit{features}). El diagrama lógico es:

\begin{center}
\small\texttt{
[AVRO/Impala]
$\;\to\;$ Ingesta (ThreadPool) $\;\to\;$ [Parquet crudo]
$\;\to\;$ Groupby icao24 + ventaneo (Dask) $\;\to\;$ [Parquet trayectorias]
$\;\to\;$ Features densas (Process+OpenMP) $\;\to\;$ [Parquet features]
$\;\to\;$ IsolationForest scoring (Dask) $\;\to\;$ [Top-k anomalías]
}
\end{center}

Los \textit{kernels} críticos (DTW entre trayectoria observada y ruta canónica por par origen-destino, distancia haversine punto a punto y suavizado Kalman simplificado) se implementan en C++ con OpenMP y se exponen a Python vía \texttt{pybind11}. El contrato es: \textit{kernel} recibe \texttt{numpy.ndarray} contiguo, libera el GIL con \texttt{py::gil\_scoped\_release}, paraleliza el bucle interno con \texttt{reduction} sobre \textit{scores} y devuelve un arreglo de tamaño fijo por trayectoria. Este aislamiento garantiza que la parte rápida no compita con el GIL.

\section*{5. Herramientas tentativas}

Python 3.11+, Dask Distributed (\textit{scheduler} + \textit{workers}), \texttt{pyarrow} para Parquet/AVRO, \texttt{numpy} y \texttt{scikit-learn} (\texttt{IsolationForest}), kernel propio en C++17 con OpenMP compilado mediante \texttt{pybind11}, almacenamiento Parquet sobre objeto compatible S3 (MinIO local para desarrollo, S3/GCS en cluster). MPI se mantiene como alternativa si el \textit{shuffle} de Dask no escala bien sobre $\geq 32$ \textit{workers}. La elección entre Dask y Spark se justifica por la afinidad nativa con el \textit{stack} numérico de Python y el menor costo de mover datos entre el \textit{scheduler} y los \textit{kernels} OpenMP.

\section*{6. Métricas de evaluación y diseño de \textit{benchmarking}}

Se miden tres dimensiones simultáneamente, separadas para no mezclar I/O con cómputo. \textbf{Rendimiento:} tiempo total por etapa, \textit{throughput} en trayectorias procesadas por minuto, $S_p = T_1/T_p$ y eficiencia $E_p = S_p/p$ para $p \in \{1, 4, 16, 64\}$ \textit{workers}, memoria pico por \textit{worker} y \textit{spill} a disco. \textbf{Calidad:} \texttt{precision@k} para $k \in \{10, 50, 200\}$ contra incidentes NTSB cruzados, \texttt{recall} sobre anomalías sintéticas inyectadas con distintas magnitudes de perturbación, estabilidad del ranking entre corridas (Spearman). \textbf{Costo:} \$/TB procesado y \$/anomalía top-50 detectada, para discutir factibilidad operacional.

El diseño experimental contempla: tres repeticiones por configuración con corrida de calentamiento, fijación de semilla en la inyección de anomalías, comparación de tamaños de chunk $\{32, 128, 512\}$ MiB para discutir granularidad, y aislamiento explícito de etapas para evitar el sesgo de medir solo el \textit{end-to-end}. Las tablas reportan media, desviación y mejor de tres, no solo el mínimo.

\section*{7. Riesgos, límites y supuestos}

El primer riesgo es el \textit{data skew}: aeropuertos como LAX o JFK aportan muchísimos más eventos que aeropuertos regionales, por lo que el \textit{groupby} por \texttt{icao24} puede generar particiones desbalanceadas y \textit{workers} ociosos. Mitigación: re-particionamiento adaptativo o \textit{salting} del \textit{key}. El segundo riesgo es la cobertura: sobre océanos abiertos OpenSky tiene gaps de cobertura, lo que produce \textit{features} sesgadas para vuelos transoceánicos; se documenta como límite y se trata con interpolación cauta o filtrado por densidad mínima. El tercero es el \textit{ground truth}: NTSB sólo reporta incidentes graves; muchas anomalías reales no aparecen como ``incidente'' formal y el modelo puede ser correcto aunque la \texttt{precision@k} cruda parezca baja. El cuarto es el costo: un año a TB-scale en cloud puede llegar a tres dígitos USD por corrida completa, por lo que el plan incluye una fase reducida sobre un mes (~10\% del costo) para iterar diseño antes de la corrida final. Otros supuestos: la API/dump de OpenSky se mantiene operativa, el formato AVRO no cambia durante el ciclo, y el cluster destino tiene compiladores con soporte OpenMP $\geq 4.5$. El alcance \textit{excluye} predicción de incidentes futuros: el proyecto detecta anomalías retrospectivas, no es un sistema de alerta en tiempo real.

\section*{8. Resultado esperado del proyecto final}

\textbf{Resultado cuantitativo objetivo}: $S_p \geq 8$ con 16 \textit{workers} ($E_p \geq 0{,}5$), procesamiento de un mes en menos de 4 horas y de un año en menos de 48 horas sobre un cluster modesto (16 \textit{workers}, 8 vCPU, 32 GiB c/u), \texttt{precision@50} $\geq 0{,}3$ contra cruce NTSB y \texttt{recall} $\geq 0{,}7$ sobre anomalías sintéticas de magnitud media. \textbf{Entregables}: código reproducible empaquetado (pipeline Dask + kernel OpenMP + Dockerfile), \textit{notebook} con análisis del top-50 de anomalías por mes, informe corto con comparación contra una referencia secuencial y discusión honesta de los casos donde el ranking falla. Si las metas cuantitativas no se alcanzan, el proyecto reporta por qué (cuello documentado con perfiles) en lugar de inflar resultados.

\end{document}
